\documentclass[conference]{IEEEtran}
\IEEEoverridecommandlockouts
% The preceding line is only needed to identify funding in the first footnote. If that is unneeded, please comment it out.
\usepackage{cite}
\usepackage{amsmath,amssymb,amsfonts}
\usepackage{algorithmic}
\usepackage{graphicx}
\usepackage{textcomp}
\usepackage{xcolor}
\usepackage[UTF8]{ctex}
\usepackage{booktabs}
\usepackage{multirow}

\def\BibTeX{{\rm B\kern-.05em{\sc i\kern-.025em b}\kern-.08em
    T\kern-.1667em\lower.7ex\hbox{E}\kern-.125emX}}
\begin{document}

\title{人脸伪造视频检测综述*\\
%{\footnotesize \textsuperscript{*}注：子标题不包含在 Xplore 中，不应使用} % 此行已注释，请勿在最终论文中保留
% \thanks{本文受 [在此处填写资助机构，如果无则删除此行] 资助。}
}

\author{\IEEEauthorblockN{1\textsuperscript{st} 陈安康}
\IEEEauthorblockA{\textit{计算机学院} \\
\textit{中山大学}\\
广州, 中国 \\
1743826979@qq.com}
% \and
% \IEEEauthorblockN{2\textsuperscript{nd} 作者姓名}
% \IEEEauthorblockA{\textit{单位部门名称} \\
% \textit{单位名称}\\
% 城市, 国家 \\
% 电子邮件地址 或 ORCID}
% \and
% \IEEEauthorblockN{3\textsuperscript{rd} 作者姓名}
% \IEEEauthorblockA{\textit{单位部门名称} \\
% \textit{单位名称}\\
% 城市, 国家 \\
% 电子邮件地址 或 ORCID}
% \and
% \IEEEauthorblockN{4\textsuperscript{th} 作者姓名}
% \IEEEauthorblockA{\textit{单位部门名称} \\
% \textit{单位名称}\\
% 城市, 国家 \\
% 电子邮件地址 或 ORCID}
% \and
% \IEEEauthorblockN{5\textsuperscript{th} 作者姓名}
% \IEEEauthorblockA{\textit{单位部门名称} \\
% \textit{单位名称}\\
% 城市, 国家 \\
% 电子邮件地址 或 ORCID}
% \and
% \IEEEauthorblockN{6\textsuperscript{th} 作者姓名}
% \IEEEauthorblockA{\textit{单位部门名称} \\
% \textit{单位名称}\\
% 城市, 国家 \\
% 电子邮件地址 或 ORCID}
}

\maketitle

\begin{abstract}
本文对人脸伪造视频（Deepfake）检测的最新进展与挑战进行了全面综述。随着人工智能技术的发展，Deepfake生成内容日益逼真，在制造虚假信息、伪造电子证据和实施网络欺诈等方面构成日益严峻的社会威胁。因此，开发鲁棒高效的Deepfake检测技术变得至关重要。
本文首先概述了Deepfake检测领域常用的评价指标。随后，系统性地分析了当前主流的检测方法，主要从三个维度展开：基于时空域分析的方法、基于内在特征和一致性线索的方法、以及基于多模态融合的方法。
本文还分析了Deepfake检测技术面临多重挑战。
并针对这些挑战，展望了未来的发展趋势，本文旨在为Deepfake检测领域的研究者提供一份全面的参考，并指引未来的研究方向。
\end{abstract}

\begin{IEEEkeywords}
人脸伪造视频, DeepFake, 检测, 评价指标, 神经网络, 多模态, 泛化能力, 鲁棒性, 数据集
\end{IEEEkeywords}

\section{引言}
随着人工智能技术的飞速发展，深度伪造（Deepfake）技术已成为数字时代一股不容忽视的力量。通过利用复杂的深度学习算法，Deepfake能够生成或篡改高度逼真的人脸图像和视频，使其真假难辨。这项技术在娱乐、艺术创作等领域展现出巨大潜力的同时，也带来了前所未有的严峻挑战。Deepfake可能被恶意用于制造虚假新闻、传播不实信息、伪造电子证据、实施网络诈骗和数字骚扰，对个人声誉、社会信任乃至国家安全构成严重威胁。

面对Deepfake技术日益成熟和广泛应用所带来的风险，有效且鲁棒的Deepfake检测技术的研发显得刻不不容缓。尽管该领域已取得了显著进展，但伪造技术的快速迭代与演进，不断对现有检测方法提出更高要求。理解和应对Deepfake的不断变化，对于维护数字媒体的真实性和社会稳定至关重要。

本文旨在对当前人脸伪造视频检测的前沿研究进行全面而系统的综述。我们将首先概述Deepfake检测领域常用的评价指标，以期为后续的模型性能评估提供统一标准。随后，本文将从时空域分析、内在特征与一致性线索分析以及多模态融合三大核心维度，对现有的Deepfake检测方法进行深入梳理和归纳。在此基础上，我们将剖析当前研究面临的主要挑战与瓶颈，例如模型的泛化能力与鲁棒性不足、特定场景检测的局限性以及伪造技术快速发展带来的检测滞后性。最后，本文将对Deepfake检测领域的未来发展趋势进行展望，旨在为研究人员提供有价值的参考，并指引未来的研究方向，以共同应对日益复杂的Deepfake威胁。

\section{评价指标}
在人脸伪造视频检测领域，准确评估模型性能至关重要。本节将介绍常用评价指标，包括准确率、精确率、召回率、F1分数、AUC曲线等。

\subsection{准确率 (Accuracy)}
准确率是最直观的评价指标，表示模型正确预测的样本数占总样本数的比例。对于分类任务而言，它衡量模型对所有类别样本的整体分类能力。其计算公式为：
$$ \text{Accuracy} = \frac{\text{TP} + \text{TN}}{\text{TP} + \text{TN} + \text{FP} + \text{FN}} $$
其中，TP (True Positive) 表示正确识别的伪造视频（真阳性），TN (True Negative) 表示正确识别的真实视频（真阴性），FP (False Positive) 表示将真实视频错误识别为伪造视频（假阳性），FN (False Negative) 表示将伪造视频错误识别为真实视频（假阴性）。

\subsection{精确率 (Precision) 和召回率 (Recall)}
精确率 (Precision) 关注模型预测为伪造视频的样本中，有多少是真正的伪造视频，反映了模型的查准能力，即“宁可不报，不能错报”。其计算公式为：
$$ \text{Precision} = \frac{\text{TP}}{\text{TP} + \text{FP}} $$
召回率 (Recall) (又称查全率或敏感度) 关注所有真正的伪造视频中，有多少被模型成功识别，反映了模型的查全能力，即“宁可多报，不能漏报”。其计算公式为：
$$ \text{Recall} = \frac{\text{TP}}{\text{TP} + \text{FN}} $$
实际应用中，精确率和召回率往往相互制约。

\subsection{F1 分数 (F1-Score)}
F1 分数是精确率和召回率的调和平均值，旨在综合衡量模型性能。当精确率和召回率出现较大偏差时，F1 分数会较低，因此它在评价不平衡数据集上的分类器性能时尤为有用。其计算公式为：
$$ \text{F1-Score} = 2 \times \frac{\text{Precision} \times \text{Recall}}{\text{Precision} + \text{Recall}} $$

\subsection{ROC 曲线与 AUC 值}
受试者工作特征曲线 (Receiver Operating Characteristic Curve, ROC) 是一种图形化工具，用于展示二分类模型在不同分类阈值下的性能。它以假阳性率 (False Positive Rate, FPR) 为横轴，真阳性率 (True Positive Rate, TPR) 为纵轴绘制。其中，$\text{FPR} = \frac{\text{FP}}{\text{FP} + \text{TN}}$，$\text{TPR} = \frac{\text{TP}}{\text{TP} + \text{FN}}$ (即召回率)。
曲线下面积 (Area Under Curve, AUC) 是对 ROC 曲线的定量衡量，其值介于 0 到 1 之间。AUC 值越高，表示模型分类性能越好。AUC 对数据集类别不平衡不敏感，因此在深度伪造检测中被广泛采用。

\subsection{其他评价指标}
除了上述常用指标，人脸伪造视频检测领域还可能采用其他评价指标，从不同角度评估模型性能：
\begin{itemize}
    \item 等错误率 (Equal Error Rate, EER)： 假阳性率 (FPR) 等于假阴性率 (FNR) 时的错误率。EER 越低，模型性能越好，表示模型在平衡两种错误时的最佳性能点。
    \item 平均精确率 (Average Precision, AP)： 常用于目标检测与信息检索，在伪造检测中可评估模型识别特定伪造类型的性能。它是精确率-召回率曲线下的面积。
    \item 检测错误权衡 (Detection Error Trade-off, DET) 曲线： 与 ROC 曲线类似，但它以 FNR 和 FPR 的对数尺度绘制，更强调低错误率区域的性能。
\item 推理速度 (Inference Speed) 和 模型大小 (Model Size)： 在实际应用中，尤其是在实时检测场景下，模型的推理速度和计算资源消耗（模型大小）亦是重要考量指标。
\end{itemize}

---

\section{人脸伪造视频检测方法}
人脸伪造视频检测方法可从不同角度分类总结。本节将从时空域、内在特征和一致性以及多模态融合三个主要维度对现有方法进行分析。

\subsection{基于时空域分析的方法}
这类方法关注视频帧内部或帧间的时空一致性、伪影或特定模式。

近年来，伪造人脸视频检测方法逐渐从单一的空间域或时间域分析，发展为融合多维度特征的综合建模策略。其中，基于时空信息的不一致性建模已成为主流方向之一。这类方法通过捕捉真实人脸视频在时空维度上的自然关联性和规律性，识别伪造视频因合成缺陷导致的帧内（空间）或帧间（时间）时空不一致性或异常。相较于早期仅关注静态图像的空间特征或视频帧间的时间变化，最新研究普遍倾向于联合建模空间与时间维度，以捕捉深度伪造视频中潜在的细微时空不一致性。此外，除传统时空特征融合外，近年研究亦引入频率域分析，通过提取图像或视频的频率特征，进一步提升模型对伪造痕迹的感知和检测准确率。

\subsubsection{深度学习驱动的时空特征学习}
以往时空域研究常采用CNN进行空间特征提取（如Xception\cite{F.-421}、ResNet\cite{K.X.-422}、InceptionNet\cite{C.Wei-423}），RNN进行时间特征提取（如LSTM\cite{S.J.-424}），形成“CNN+RNN”架构。随着深度学习技术发展，研究者们逐渐引入更先进的神经网络架构，如Transformer等，以实现时空特征的更深层次挖掘。

研究者常以 Transformer 替换传统框架中的 RNN，构建“CNN+Transformer”框架。针对低层空间特征学习与时间信息中时序变化捕捉的关键技术难题，Weirong Sun 等人（2023）创新性提出卷积-Transformer 深度伪造检测模型（ConTrans-Detect），构建双模块协同架构，其中多尺度CNN模块借助3D Inception块高效提取视频多尺度低层特征，为空间特征表征奠定基础；多分支 Transformer 模块由多流 Transformer 层构成，各层输入差异化的时间分辨率与空间特征维度，精准感知视频复杂运动变化，完成时序特征深度建模。实验结果表明，在 DFDC 数据集上斩获 0.929 的 AUC 值与 0.920 的 F1 分数，显著超越多项先进检测方法，展现出卓越的深度伪造检测性能\cite{WY406}。
Yu等人（2023）提出一种多时空视角 Transformer（MSVT），包含局部时空视角（LSV）和全局时空视角（GSV）模块。研究者使用ResNet-50提取每帧特征。不同于常规稀疏采样单帧构建输入序列，其采用局部连续时间视角捕捉关键动态不一致性。每组提取的帧特征送入时间 Transformer，生成组级别时空特征。所有帧特征送入时间Transformer，并通过特征融合模块建立全局时空视角（GSV）。最后，获取的局部和全局特征通过全局 - 局部 Transformer（GLT）整合，挖掘更微妙、全面的特征。在FF++等数据集上，准确率达 98.48\%，优于其他先进检测方法\cite{YuNi-129}。
现有基于时空特征的深度伪造检测模型在学习机制本质的可解释性层面存在显著局限。针对此关键问题，Zhao等（2023）构建可解释时空视频 Transformer 模型（ISTVT）。其核心创新在于设计分解式时空自注意力模块与自减法机制，差异化捕捉空间伪造痕迹特征与时间序列不一致性，显著提升模型鲁棒判别能力。该模型基于分解式架构设计，使用 Xception 网络入门流部分作为特征提取器，专注于捕获低级纹理信息，因伪造细微痕迹常存在于此。为解决帧间不一致性，分解式时空自注意力模块将空间和时间维度注意力分开处理，更精确识别单帧内部伪影（如模糊边缘）和帧序列间动态异常（如不连贯的嘴唇运动）。在此基础上，自减法机制通过计算相邻时间步的特征残差，促使模型聚焦关键帧间失真信息，忽略冗余背景变化。更重要的是，ISTVT 借助相关性传播算法（一种输入像素级重要性映射技术），对时空维度关键判别区域进行可视化映射，建立模型决策逻辑与特征空间的可解释关联，为理解深度伪造检测模型决策机制提供直观认知路径。在 FaceForensics++ 等多源异构大规模数据集上的实证研究表明，所提方法在同数据集内检测任务与跨数据集泛化场景中均展现领先性能，充分验证其有效性与鲁棒性。\cite{ZhaoWang-68}。

随着深度伪造视频检测研究深入，研究者亦依据对问题的不同理解，提出多种新颖的时空特征学习方法。
Lu 等人（2024）将伪造人脸检测视为一种特殊细粒度图像分类任务，提出基于长距离注意力机制的方法，从空间和时间两个维度全局建模伪造痕迹。该方法有效整合全局上下文信息，同时增强局部统计特征表达能力。在多个公开数据集上的实验结果表明，其检测精度表现优异\cite{LuLiu-167}。
Pang 等人（2023）针对深度伪造视频中常见的动态时空不一致性问题，提出新颖的多速率激励网络（Multi-rate Excitation Network, MRE-Net）。该方法从多速率角度激励特征，通过创新两分图组采样策略（BGS）、负责瞬时空间和短期时间不一致性的瞬时不一致性激励模块（MIE），以及编码长时时间不一致性的长时不一致性激励模块（LIE），有效增强模型对时空不一致性伪造痕迹的感知能力。实验在 FF++ 和 DFDC 等主流数据集上进行，结果显示 MRE-Net 在检测准确率上显著优于 MesoNet、Xception 等主流伪造检测方法，展示其在视频伪造检测任务中的强大性能\cite{PangZhang-372}。

\subsubsection{融合频率域等信息以增强时空分析}
近年来，随着人脸伪造检测研究深入，时空域特征分析逐步与其它维度信息（尤其是频率域）有机融合，通过多维度信息互补，显著提升模型对复杂伪造模式的泛化能力和鲁棒性。
Yue 等人（2024）分析深度伪造视频中全局与局部帧间动态不一致性，提出局部区域频率引导的动态不一致网络（Local-region Frequency-Guided Dynamic Inconsistency Network, LFGDIN）。该网络由全局时空网络（Global Spatio-Temporal Network, GSTN）与局部区域频率引导模块（Local Region Frequency-Guided Module, LRFGM）两部分组成，其中 GSTN 建模面部全局动态特征，LRFGM 则专注于眼部与口部的频率动态信息提取。LRFGM 通过频率域信息引导 GSTN 关注关键局部区域，增强模型对动态不一致性的建模能力。在 FF++、DFDC 和 Celeb-DF 三个公开数据集上，LFGDIN 在检测未知伪造方法视频时，相较多种先进方法展现更优检测性能\cite{YueChen-66}。
为提升模型鲁棒性，尤其是在复杂场景下对伪造视频的判别能力，Emara 和 Elagamy（2024）提出双流深度伪造视频检测框架 DeepStream-X，通过频率特征与时空特征高效融合增强检测性能。该框架构建双路径特征提取机制：其一利用频率分析挖掘单帧内细微异常变化，其二通过密集光流运动估计捕捉跨帧间运动不一致性，形成时空-频率特征互补表征。融合特征经引入注意力机制的改进 Xception 网络层级化学习，有效提升判别性特征提取能力。实验验证表明，在 Celeb-DF v2 数据集上，DeepStream-X 框架实现了 95.24\% 的检测准确率与 99.45\% 的 AUC 值，显著优于同类方法，展现出复杂场景下对深度伪造视频的精准判别能力与强鲁棒性\cite{NM414}。

研究者还考虑到引入频率特征能够优化传统神经网络模型，因此着手利用频率域特征改进现有卷积神经网络结构。为优化卷积内部特征表示，构建适用于深度伪造检测的新型主干网络结构，Guo等人（2024）提出一种新型空频交互卷积（Space-Frequency Interactive Convolution, SFIConv）。该网络使用多通道约束可分离卷积（Multichannel Constrained Separated Convolution, MCSConv）作为核心组件，提取高频特征，通过生成、交互、融合三阶段实现空频特征联合学习。SFIConv作为通用网络结构，可替换任意主干网络中的标准卷积。大量实验表明，嵌入SFIConv的主干网络可显著提升深度伪造检测准确率，且空频交互机制有助于捕获通用伪影特征，使其在跨数据集评估中表现更优\cite{GuoJia-191}。

\subsubsection{考虑干扰因素的时空分析}
除多维度信息互补外，亦有研究引入对其他维度潜在干扰信息的考虑，通过去噪机制，显著提升模型对复杂伪造模式的判别能力。
Yin等人（2023）考虑面部运动带来的干扰，研究长短距离帧间运动，提出动态差异学习方法，区分人脸伪造引起的帧间差异与面部自然运动引起的帧间差异，实现深度伪造视频检测中精准的时空不一致性建模。该方法包含动态细粒度差异捕捉模块（DFDC）与多尺度时空聚合模块（MSA），前者通过自注意力机制与去噪操作，消除面部运动导致的干扰差异并生成长距离差异注意力图；后者则通过聚合多方向、多尺度的时间信息强化时空不一致性表征。这两个模块可无缝集成至现有二维卷积神经网络，构建动态时空不一致性捕捉网络。大量实验表明，在不同基准数据集上，算法检测性能显著优于多项主流方法，展现稳健领先优势\cite{YinLu-279}。

% 在 \subsection{基于时空域分析的方法} 的文字描述后，具体方法介绍前插入此表格
\begin{table*}[htbp] % table* 用于跨双栏，htbp 表示表格可以放置的位置
\centering
\caption{基于时空域分析的Deepfake检测方法总结} % 表格标题
\label{tab:spatial_temporal_methods} % 表格标签，用于交叉引用
\resizebox{\textwidth}{!}{ % resizebox 用于调整表格宽度以适应页面，如果表格过宽
\begin{tabular}{@{}lccccl@{}} % @{} 表示不留左右间距，l: 左对齐, c: 居中, r: 右对齐
\toprule
\textbf{方法类别} & \textbf{代表模型/算法} & \textbf{核心思想} & \textbf{关键优势} & \textbf{主要数据集} & \textbf{实验结果} \\
\midrule
\multirow{3}{*}{深度学习驱动的时空特征学习}
& ConTrans-Detect \cite{WY406} & CNN+Transformer 双模块协同，多尺度CNN提取低层特征，多分支Transformer感知复杂运动 & 高效提取多尺度时空特征，性能卓越 & DFDC & AUC 0.929, F1 0.920 \\
& MSVT \cite{YuNi-129} & 多时空视角Transformer（LSV+GSV），连续时间视角捕捉动态不一致性 & 捕捉局部与全局细微特征，高准确率 & FF++ & 准确率 98.48\% \\
& ISTVT \cite{ZhaoWang-68} & 可解释时空视频Transformer，分解式自注意力与自减法机制 & 提升模型鲁棒判别能力，提供决策可解释性 & FF++ & 同数据集与跨数据集泛化均表现优异\\
\midrule
\multirow{2}{*}{融合频率域等信息}
& LFGDIN \cite{YueChen-66} & 全局时空网络+局部区域频率引导模块，关注眼部/口部频率动态 & 增强动态不一致性建模能力，泛化性好 & FF++, DFDC, Celeb-DF & 优于多种先进方法（未知伪造方法检测） \\
& DeepStream-X \cite{NM414} & 双流框架，频率特征与时空特征高效融合，引入注意力机制 & 提升复杂场景鲁棒性，高准确率AUC & Celeb-DF v2 & 准确率 95.24\%, AUC 99.45\% \\
& SFIConv \cite{GuoJia-191} & 新型空频交互卷积，替换标准卷积，提取高频特征 & 提升检测准确率，捕获通用伪影特征，跨数据集表现优 & FF++, Celeb-DF, DFFD & 显著提升主干网络检测准确率 \\
\midrule
\multirow{1}{*}{考虑干扰因素的时空分析}
& 动态差异学习 \cite{YinLu-279} & 区分伪造与自然面部运动差异，动态细粒度差异捕捉+多尺度时空聚合 & 消除面部运动干扰，精准建模时空不一致性 & FF++, Celeb-DF & 显著优于多项主流方法（如F3Net、Xception、Mesonet等） \\
\bottomrule
\end{tabular}
} % resizebox 结束
\end{table*}
---

\subsection{基于内在特征和一致性线索的方法}
随着人脸伪造技术不断进步，传统基于外观特征的检测方法逐渐暴露出局限性。研究者开始关注人脸视频中难以被伪造技术精确复制的内在特征和一致性线索。

\subsubsection{基于生理信号的方法}
心率、呼吸等生理信号是真实人脸固有生物特征，伪造视频难以精确复制其微弱规律信号。
早期生理信号研究主要集中于面部微运动，特别是眼球运动/眨眼模式\cite{Y.M.-426}。目前，面部生理信号研究日益多样化。
针对压缩视频产生的额外伪影导致的检测难题，Liao等人（2023）提出基于面部肌肉运动（FAMM）的框架。FAMM包含面部地标提取、面部肌肉运动特征构建及预测概率融合三个模块。该框架首先从连续视频帧中精确定位人脸，提取关键面部地标。随后，研究者从两个角度进行特征转换，构建面部肌肉运动特征：一方面，通过计算相邻帧间面部距离和角度变化的差分特征，训练门控循环单元（GRU）分类器捕捉不自然面部运动信息；另一方面，为应对压缩噪声对时间维度建模的挑战，他们计算面部肌肉运动特征的时间序列，训练支持向量机（SVM）分类器。最终，为增强检测鲁棒性，该研究利用德姆斯特-沙弗理论（Dempster-Shafer theory）融合GRU和SVM分类器输出，提供最终检测结果。德姆斯特-沙弗理论（Dempster-Shafer Theory, DST）\cite{Shafer1976}，亦称证据理论，是一个处理不确定性和不精确信息的数学框架。它允许将信任分配给命题子集，而非单个命题，并能有效融合来自多个证据源的信息，即便这些证据存在不确定或冲突。该研究通过对压缩视频的互信息、压缩过程及面部地标的详尽理论分析，证实FAMM在压缩环境下仍能有效构建面部肌肉运动特征，捕捉真实与伪造视频间的固有差异，从而在检测压缩深度伪造视频方面超越现有先进方法，并在真实社交网络环境中保持出色检测鲁棒性\cite{LiaoWang-137}。
Peng等人（2024）关注视频中注视特征研究。通过注视方向分析，研究者发现伪造与真实视频在注视方向模式分布上存在差异。相比真实视频，伪造视频具有更分散的注视分布。针对此特征，研究者提出名为DFGaze的新型深度伪造注视分析方法，探索视频面部伪造检测中的时空注视不一致性。DFGaze 使用注视分析模型（GAM）分析面部视频帧的注视特征，通过时空特征聚合器，结合注视特征实现真实性分类。为更有效捕捉真实与伪造视频间的细微差异，模型进一步整合纹理分析模型（TAM）和属性分析模型（AAM），增强时空特征表达能力，更深入挖掘视频中真实性线索。大量实验表明，此方法实现优秀性能\cite{PengMiao-144}。

随着伪造技术发展，早期基于某些生理信号的简单检测方法，其准确率和鲁棒性显著下降。研究人员开始寻求更深层、更难复制的生理痕迹。目前，光电容积脉搏波描记术（PPG）成为提取生物特征的主流手段。传统PPG通过LED光探测皮肤下血液容积的周期性变化，产生反映血流动态的PPG波形，包含心率、血氧饱和度等丰富生理信息。在伪造视频检测中，该技术并非使用外部LED，而是利用视频本身记录的环境光在人脸皮肤上产生的微弱颜色和亮度波动。尽管肉眼难以察觉，通过信号处理技术，可从这些像素变化中提取远程PPG（rPPG\cite{Ming-ZherPoh-425}）类生理信号。由于当前先进伪造技术难以精确模拟或复制这些符合生理规律的细微波动，它们成为区分真实与合成视频的重要隐式特征。
Umur Aybars Ciftci 等人（2020）认为真实人像视频中固有的微弱生理信号在伪造视频中往往无法在细节上或时间上被准确复制，可作为识别假视频的有效隐性特征。他们基于绿色通道PPG（G-PPG或G*）和基于色度的PPG（C-PPG或C*），通过对人像视频进行特定信号变换，在成对分离问题上实现了高达99.39\%的准确率，初步验证了生物信号区分真实与伪造内容的可行性。进一步地，基于这些信号变换及其相应特征集，研究者构建通用分类器，并最终引入卷积神经网络（CNN），构建基于生物信号的伪造人像视频检测器（FakeCatcher）。该研究还构建并发布“野外”伪造人像视频数据集。在 Face Forensics 等多个基准数据集及其自建数据集上的评估结果显示，该方法（FakeCatcher）检测准确率表现优异，分别达到 96\%、94.65\%、91.50\% 和 91.07\%。该方法在检测率上显著优于现有基线，且其性能独立于伪造内容来源、生成器或特定属性，展现良好泛化能力\cite{CiftciDemir-93}。

随着rPPG研究深入，当前研究亦尝试将PPG信号与其他分析技术、多维度特征融合，以更鲁棒方式捕捉生理信号不一致性。
Sun等人（2022）在深入分析现有深度伪造检测方法局限性基础上，从面部像素变化的时空表示角度切入，提出名为FakeTransformer的新颖方法。该研究核心假设在于，现有生成模型难以完美还原面部毛细血管血容量变化引起的生理性颜色变化，导致生成内容中生理信息时间序列不连续；且真实人脸毛细血管固定分布模式导致不同面部区域周期性颜色信号存在相位差，而伪造面部则缺乏这种空间一致性。为克服传统rPPG信号提取方法易受光照、肤色或运动干扰的缺点，并避免线性降维可能消除关键伪造误差的问题，该方法首先在多个高斯尺度上应用欧拉视频放大（EVM）技术，放大面部血容量变化引起的微弱生理信号，凸显面部因血流变化产生的生理性律动。随后，将经EVM处理的视频与原始视频一同转换为多尺度欧拉放大时空图（MEMSTmap）。MEMSTmap能有效表示不同尺度上随时间变化的生理增强序列，捕捉肉眼难以察觉的生理波动。为捕捉MEMSTmap中复杂细微的时空模式，研究者创新性地将这些图重塑为列单位的帧块，并输入至视觉变换器（Vision Transformer, ViT）。ViT在此被巧妙用于编码多变量时间序列，利用强大全局自注意力机制学习帧级别的鲁棒时空描，揭示视频中的空间和时间伪影。最终，通过整合特征嵌入，模型输出判断视频真伪的概率。在FaceForensics++和深度伪造检测数据集上的验证结果表明，该模型在伪造检测方面表现出色，并展示优异的跨数据域泛化能力，充分证明了从放大生理变化中提取时空生理信息进行检测的有效性\cite{Y.Z.-402}。
Yang等人（2023）考虑面部特征的时空相关性。通过自回归（AR）系数和光电容积脉搏波描记术（PPG）提取的特征作为时空域取证基础，并依据这些特征生成可靠的指纹信息。PPG 可通过视频估算心率，AR 系数则反映像素间相关性及生成假脸过程中由上采样引起的平滑空间痕迹。研究者利用获取的指纹信息，训练两个基于ACBlock的DenseNets作为分类器。该模型在多个深度伪造数据集上提供最先进性能，并展现更好的泛化能力\cite{JY-412}。

\subsubsection{基于物理光影与图像结构一致性特征}
除生理信号外，物理光影和图像结构等方面的不一致性亦是伪造视频难以精确模拟的内在一致性线索。研究者对这些细节特征的深入挖掘，能有效提升伪造视频检测的准确性和鲁棒性。
通过分析身体各部位（如眼睛、鼻子、脸颊等）的反射及环境因素，Mohzary等人（2024）提出新颖的反射和环境感知深度伪造（READFake）检测技术。READFake 为分层结构，组件划分为训练数据标注（TDA）、反射和环境因素检测（REFD）以及特征提取、嵌入和分类（FEEC）三个模块。READFake 技术通过执行 FIEP（反射和环境感知）特征提取和镜面高光检测两大任务。其中 FIEP 参数（室内/室外、光强度、光水平）通过 MobileNet-V2 模型训练和 LAB 颜色空间分析获取；CSH 和 BSH（镜面高光）则分别通过 MobileNetV2-SSDLite 训练和 HSV 颜色空间分析识别，旨在验证深度伪造内容与环境反射及光照缺乏协调性。FEEC 模块利用从 REFD 模块提取的镜面高光特征和 FIEP，执行深度层次特征提取、相似性度量、反射和环境因素嵌入以及分类四个主要功能。实证结果表明，READFake 在检测复杂深度伪造图像方面达到 99.00\% 的高准确率\cite{MohzaryBasunduwah-365}。
Deng等人（2022）关注现有伪造算法在生成视频时面部边缘留下的合成伪造痕迹，提出一种新伪造视频检测方法。该方法首先从视频帧中找到面部边缘，提取面部边缘带作为深度学习输入，基于EfficientNet-B3训练，实现对深度伪造视频的有效检测。实验表明，该方法在 Face-Forensics++ 数据集所有四种伪造方法上均实现超过 99.8\% 的 AUC 值\cite{Z.B.-410}。


% 在 \subsection{基于内在特征和一致性线索的方法} 的文字描述后，具体方法介绍前插入此表格
\begin{table*}[htbp]
\centering
\caption{基于内在特征和一致性线索的Deepfake检测方法总结}
\label{tab:intrinsic_clues_methods}
\resizebox{\textwidth}{!}{
\begin{tabular}{@{}lccccl@{}}
\toprule
\textbf{方法类别} & \textbf{代表模型/算法} & \textbf{核心思想} & \textbf{关键优势} & \textbf{主要数据集} & \textbf{实验结果} \\
\midrule
\multirow{4}{*}{基于生理信号的方法}
& FAMM \cite{LiaoWang-137} & 基于面部肌肉运动特征，融合GRU和SVM分类器，利用DST处理不确定性 & 有效应对压缩视频，鲁棒性强 & FF++, Celeb-DF & 超越传统方法（压缩视频检测） \\
& DFGaze \cite{PengMiao-144} & 注视特征分析，结合纹理与属性分析增强时空特征 & 捕捉注视模式差异，全面挖掘真实性线索 & FF++, WildDeepfake, Celeb-DF, DFDC & 在数据集内评估和跨数据集评估结果超越传统方法 \\
& FakeCatcher \cite{CiftciDemir-93} & 基于绿色通道/色度PPG信号，构建CNN检测器 & 泛化能力强，独立于伪造内容来源 & FF++, Celeb-DF, 自建数据集 & 准确率 96\% (FF++), 91.07\% (自建) \\
& FakeTransformer \cite{Y.Z.-402} & 放大生理信号并转为MEMSTmap，ViT编码多变量时间序列 & 捕捉微弱生理波动，跨数据域泛化能力优异 & FF++, DFDC & 表现出色，优异跨数据域泛化 \\
& PPG+AR+DenseNet \cite{JY-412} & PPG与AR系数提取指纹信息，基于ACBlock的DenseNets分类 & 性能先进，泛化能力好 & 多种Deepfake数据集 & 提供先进性能 \\
\midrule
\multirow{2}{*}{基于物理光影与图像结构}
& READFake \cite{MohzaryBasunduwah-365} & 分析面部反射及环境光影因素，FIEP与镜面高光检测 & 高准确率，验证伪造内容与环境光照缺乏协调性 & 自建数据集 & 准确率 99.00\% \\
& 面部边缘带检测 \cite{Z.B.-410} & 提取面部边缘带作为深度学习输入，基于EfficientNet-B3训练 & 高AUC值，有效识别合成伪造痕迹 & FF++ & AUC 超 99.8\% (所有伪造方法) \\
\bottomrule
\end{tabular}
}
\end{table*}
---

\subsection{基于多模态融合的方法}
这类方法结合视觉、音频甚至文本等不同模态信息，以提高检测的鲁棒性和准确性。目前在伪造人脸检测领域，多模态融合主要集中于视觉和音频模态的协同分析。由于深度伪造技术能操纵视频的视觉与音频信号，常导致两者间不一致，为多模态检测提供契机。
Shahzad等（2022）提出新颖的基于唇读的多模态深度伪造检测方法“Lip Sync Matters”。该方法通过分析视频中唇形序列与音频模型 Wav2lip 生成的合成唇形序列间的高层语义不匹配来识别伪造视频\cite{S.A.-427}。但面对不包含音频操控的伪造视频区分能力较弱。
Oorloff等人（2024）提出了一种两阶段的跨模态学习方法——音视频特征融合（AVFF）。第一阶段通过在真实视频上进行自监督学习，利用对比学习和自编码目标，并结合音频-视觉互补掩蔽和特征融合策略，来学习丰富的跨模态表示。第二阶段则通过在真实和虚假视频上进行监督学习，对已学习的表示进行微调以进行深度伪造分类。该方法在FakeAVCeleb数据集\cite{khalid2021fakeavceleb}上取得了98.6\%的准确率和99.1\%的AUC。但不能检测单模态的视频\cite{OorloffKoppisetti-290}。

除视觉和音频模态协同分析外，研究者亦开始关注如何将其他模态（如文本）与视觉模态融合，以提升伪造视频检测准确性。Song等人（2024）提出创新多模态检测（MM-Det）算法。该方法融入文本（指令）作为引导 LMM 理解和推理的关键元素。利用大语言多模态模型（LMM）强大感知和综合能力，通过从 LMM 多模态空间生成多模态伪造表示（MMFR），显著增强其检测扩散模型生成视频（特别是未见伪造内容）的能力。此外，MM-Det 框架结合帧内和帧间注意力（IAFA）机制在时空域进行特征增强，并通过动态融合策略优化伪造表示组合。为支持扩散视频检测研究，作者构建全面扩散视频取证（DVF）数据集，涵盖多种主流扩散生成方法\cite{SongGuo-428}。

当前深度伪造技术向视觉-音频协同伪造发展，模态间不一致性变得更难察觉，给现有方法带来挑战。随着多模态在伪造人脸检测领域研究深入，模态对齐（解决并弥合不同模态间表示差距）成为研究者关注重点。针对此问题，Yu等人（2024）提出预测性视觉-音频对齐自监督（PVASS-MDD）框架，以更有效提取VA不一致性。该框架包含两个阶段：首先是PVASS辅助阶段，通过设计三流网络关联增强视觉视图与音频线索，并通过新颖的跨模态预测对齐模块消除模态差距，从而在真实视频中学习固有的VA对应关系。其次是MDD阶段，利用辅助损失并结合冻结的PVASS网络对齐真实视频VA特征，以指导多模态检测器捕捉更微妙的VA不一致。实验证明，此类方法在现有数据集上能显著提升检测性能，但在面对光照不足和极端面部角度的情况下，难以提取有价值的信息，导致模型能力不足\cite{YuLiu-149}.
Liu等人（2024）提出利用多模态对比学习（MCL）的新型深度伪造检测方法，旨在更好地探索单模态和跨模态伪造线索。该方法通过跨模态对比学习策略，从多模态信息中学习组合嵌入，拉近单模态与多模态表示差距，并挖掘伪造痕迹。具体而言，它首先分别提取音频、视频帧及整体视频三种模态特征，随后将音频和帧特征与视频特征组合，形成跨模态表示，并通过对比学习缩小跨模态特征与单模态特征间的差距。这种MCL方法能联合学习包含单模态和跨模态伪造线索的更有效表示。此外，为增强视频网络对帧内伪造线索的挖掘能力，还引入帧知识蒸馏，且不增加额外计算开销。最后，通过基于噪声的特征增强（NFA）模块自适应扰动视听特征，进一步提升泛化性能。实验结果表明，该框架显著优于现有最先进方法\cite{LiuYu-76}.

% 在 \subsection{基于多模态融合的方法} 的文字描述后，具体方法介绍前插入此表格
\begin{table*}[htbp]
\centering
\caption{基于多模态融合的Deepfake检测方法总结}
\label{tab:multimodal_methods}
\resizebox{\textwidth}{!}{
\begin{tabular}{@{}lccccl@{}}
\toprule
\textbf{方法类别} & \textbf{代表模型/算法} & \textbf{核心思想} & \textbf{关键优势} & \textbf{主要数据集} & \textbf{实验结果} \\
\midrule
\multirow{1}{*}{视觉-音频协同分析}
& Lip Sync Matters \cite{S.A.-427} & 基于唇读，分析唇形序列与合成音频的高层语义不匹配 & 识别视觉与音频不一致性，直观有效 & FakeAVCeleb\cite{khalid2021fakeavceleb} & 在 Faceswap wav2lip、Fsgan wav2lip、Wav2lip、Test-set-1 和 Test-set-2 的情况下，实现了 98\%、97\%、96\%、94\% 和 94\% 的准确率。在 Faceswap、Fsgan 和 RTVC 测试集上的表现略差\\
& AVFF \cite{OorloffKoppisetti-290} & 两阶段跨模态学习，音频-视觉互补掩蔽与特征融合 & 学习丰富的跨模态表示，检测性能优异 & FakeAVCeleb\cite{khalid2021fakeavceleb} & 准确率 98.6\%, AUC 99.1\%，但不能检测单模态的视频\\
& PVASS-MDD \cite{YuLiu-149} & 预测性视觉-音频对齐自监督，消除模态差距学习VA对应关系 & 有效提取VA不一致性，提升检测性能 & DFDC、FakeAVCeleb\cite{khalid2021fakeavceleb}、DefakeAVMiT\cite{yang2023avoid} & 数据集内测试优于RealForensics、AVoiD-DF、AVAD等模型，参数量相比于 AVAD增加了 8.1M，在光照不足和极端面部角度情况下检测表现不佳 \\
& MCL \cite{LiuYu-76} & 多模态对比学习，探索单模态和跨模态伪造线索 & 联合学习有效表示，增强泛化性能 & DF-TIMIT、DFDC、FakeAVCeleb\cite{khalid2021fakeavceleb}、KoDF\cite{Kwon2021KoDF} & 数据集内和跨数据集泛化测试中均显著优于Joint Audio-Vsual等先进方法 \\
\midrule
\multirow{1}{*}{文本-视觉融合}
& MM-Det \cite{SongGuo-428} & 融入文本（指令）引导LMM理解推理，生成多模态伪造表示 & 增强对扩散模型生成视频的检测能力，特设DVF数据集 & 自建库 & AUC 95.7 \\
\bottomrule
\end{tabular}
}
\end{table*}
---

\section{挑战与瓶颈}
人脸伪造视频检测技术正快速演进，面临持续挑战与广阔创新空间。本节将探讨亟待解决的关键问题。

模型泛化能力与鲁棒性不足。
目前，人脸伪造视频检测模型往往在特定数据集上表现出色，但其泛化能力和鲁棒性在面对不同数据集或真实应用场景时显得捉襟见肘，实用价值受限。Nadimpalli等人（2022）测试了伪造人脸检测领域主流的卷积神经网络模型在跨数据集时的表现，结果显示，ResNet-50、XceptionNet、EfficientNet V2-L等模型在跨数据集测试时准确率大幅下降。这表明当前模型在不同数据集间的泛化能力较差，难以适应实际应用中的多样化需求\cite{A.A.-429}。
与此同时，目前的大多数模型针对压缩视频的检测能力有限。由于视频压缩会引入噪声和失真，这导致其产生的伪影与深度伪造视频中的篡改伪影相互叠加，从而限制了深度学习模型对关键细节的利用。虽然目前的研究已开始关注压缩视频的检测问题（如Liao等人提出的FAMM模型\cite{LiaoWang-137}，F3-Net模型\cite{M.Z.-432}），但有研究指出这些模型实际应用中仍面临复杂性或通用性挑战\cite{ChenLiao-250}。
此外，数据集的纯洁性也会成为影响模型泛化能力的关键因素。目前的大部分的模型训练效果依赖于数据集的质量，但在真实的Deepfake检测应用场景中，数据标注者可能会误标数据，或者恶意攻击者会故意在训练数据中注入带有错误标签的样本。一旦模型在这样的数据集上训练，其泛化能力和鲁棒性将受到严重影响。研究表明，数据集中的标签噪声会显著降低模型的检测性能\cite{QiaoXie-152}。

针对特定场景检测的研究不足。现有研究大多聚焦于通用型伪造视频检测，而针对特定应用场景（如口罩遮挡、多人脸检测等）的定制化检测方法仍相对较少，难以满足多样化的实际需求。尽管研究者已经开始在此领域进行数据集构建和相关研究，如在口罩遮挡场景下的深度伪造口罩数据集（DFFMD）\cite{AlnaimAlmutairi-36}，但整体上仍缺乏足够的研究和应用。

伪造技术演进与检测方法滞后。人脸伪造技术的飞速发展，检测方法的发展速度与其不匹配，部分流行检测技术或将逐渐失效。最新研究中，Seibold等人（2025）实验证明现代高质量 Deepfake 视频能从驱动视频中继承并保留真实 rPPG 信号。这推翻了“深度伪造生成过程中，微妙心跳相关信号会丢失”的假设，表明仅依赖 rPPG 信号存在性检测伪造的早期方法或不再有效\cite{Seibold2025HighQuality}。
也有研究指出，随着伪造技术不断进步，单一模态或特征的检测方法逐渐失效\cite{Amirabadi2024DeepFakeReview}。

---

\section{未来发展趋势}
人脸伪造视频检测技术仍在快速发展，面临持续挑战。本节将探讨未来研究方向和发展趋势。

数据集的构建和多样性是伪造人脸检测研究的基础，未来将继续扩展数据集规模和多样性，以支持更全面的模型训练和评估。
当前Deepfake数据集并不全面，某些研究领域缺少专门的训练数据集，某些场景的数据集仍有空白，在未来，这些领域的数据集会被完善补充，为模型训练提供更丰富的样本和特征。研究者们将继续扩展数据集的规模和多样性，以支持更全面的模型训练和评估。
除了扩充数据集的广度和深度，另一前沿方向是提升模型应对数据本身缺陷的能力。有研究 \cite{QiaoXie-152} 开始从消除噪声标签对模型训练的负面影响入手，目标是在训练数据不完美的情况下，也能训练出高性能、高鲁棒性的Deepfake检测器。这反映了领域正从单纯依赖大规模“干净”数据，转向提升模型在真实、复杂数据环境下的适应性。

多维度特征与多模态融合已成为未来伪造人脸识别检测的重要趋势。随着技术的发展和研究的深入，越来越多的研究者尝试将不同维度，不同模态结合，期望多特征，多模态互补融合可以提升模型的鲁棒性和泛化能力。目前最常见的融合是时空域和频域的融合、时空域和生理信号的融合、视觉和音频的多模态融合。未来的研究会更加关注不同特征模态之间的对齐和融合方式（Liu等人（2024））。未来随着技术发展和研究深入，越来越多的模态和特征会被综合考虑。

先进学习模型和学习方法在伪造人脸检测中展现出强大潜力，未来会有更多的先进学习模型在该领域得到应用。先进的学习模型（如VIT）和学习方法（自监督，迁移学习）开始逐渐在伪造人脸检测领域得到应用，例如，Khormali等人（2024）使用自监督方法来进行深度伪造检测\cite{KhormaliYuan-366}，Calvert等人（2024）利用迁移学习提升Deepfake检测模型在特定文化区域的性能\cite{CalvertNg-350}。未来研究将继续挖掘这些模型和方法在该领域的潜力。

对特定场景的定制化检测方法（压缩视频、口罩）将成为未来研究重点。随着伪造技术的不断演进，针对特定伪造类型或应用场景（如口罩遮挡、多人脸识别、基于不同人种的伪造人脸检测）的定制化检测方法将更具实用价值。
---

\section*{致谢}
% 您的实际致谢内容应在此处填写，例如：
% 本研究得到 [在此处填写资助机构名称] 的资助。
% 感谢 [在此处填写感谢的对象，如同事、家人等] 的支持和帮助。

\bibliographystyle{IEEEtran}
\bibliography{shiKongYu}
\end{document}